\subsection{An evaluated wild class-three arithmetic quotient}

\begin{statement}{E.78}{Proved}{Actual wild operators, both lifting images, and complete class-three fibres}
Put $K=\mathbb Q(\zeta_3)$, $k=K_3=\mathbb Q_3(\zeta_3)$,
$\lambda=1-\zeta_3$, and $a=\zeta_3\lambda^{-3}$.
Retain E.54's actual map
$h=(x+y)/(x+\zeta_3y)$ from the nine-punctured cubic,
its source tangent $s=z/x$ of value one, and its induced
surjection $N\twoheadrightarrow F$ onto the free pro-three
group on the positive target generators $x,y$.
Write
\[
 C=F/\gamma_4F,\qquad C_2=F/\gamma_3F,\qquad
 C_1=F^{\rm ab},\qquad
 c=[x,y],\quad d=[x,c],\quad e=[y,c].
 \tag{E78.1}
\]
Commutators are $[u,v]=uvu^{-1}v^{-1}$ and Hall words are
$x^A y^B c^n d^r e^s$.

Fix positive coherent roots $3_n=3^{1/3^n}$ and
$\zeta_{3^{n+1}}$. There is a unique compatible choice
of $\lambda_n$ satisfying
\[
 \lambda_n^{3^n}=\lambda,\quad
 \lambda_n^2=-3_n\zeta_{3^{n+1}},\quad
 \lambda_{n+1}^3=\lambda_n.
 \qquad
 a_n=\zeta_{3^{n+1}}\lambda_n^{-3}.
 \tag{E78.2}
\]
Use the punctured-disc comparison determined by these
roots in E.54. If $\kappa=\kappa_3(3)$, its exact path
character is
\[
 \kappa_a=-\frac{\chi-1}{6}-\frac32\kappa.
 \tag{E78.3}
\]
Here and below every divided expression is integral in
$\mathbb Z_3$ on its stated domain.

The source marking is the following actual transport
of E.37's Frobenii. Start with
$\tau_0,\tau'_0=c_{\rm ar}\tau_0c_{\rm ar}$,
where $\chi(\tau_0)=7$ and
$t_0=\kappa(\tau_0)\in1+3\mathbb Z_3$.
The symbol $c_{\rm ar}$ denotes arithmetic complex
conjugation, distinct from the geometric Hall word $c$.
Set
\[
 \begin{gathered}
 \rho_0=\tau'_0\tau_0^{-1},\qquad
 h_0=\rho_0^{(t_0^{-1}-1)/12},\\
 \tau=h_0\tau_0h_0^{-1},\qquad
 \tau'=c_{\rm ar}\tau c_{\rm ar}.
 \end{gathered}
 \tag{E78.4}
\]
These are the E.49 free basis of the same actual group
$P=G_{K,\{3\}}^{(3)}$, and
\[
 (\chi(\tau),\kappa_a(\tau))=(7,-5/2),
 \qquad
 (\chi(\tau'),\kappa_a(\tau'))=(7,1/2).
 \tag{E78.5}
\]
The map from the old to the new free basis and the
root/path conjugations are retained in the proof.

For $\mu\in\{-5/2,1/2\}$ define $T_\mu$ on $C$ by
\[
 \begin{aligned}
 T_\mu(x)&=x^7,\\
 T_\mu(y)&=y^7c^{7\mu}
                d^{7\mu(\mu-1)/2}e^{14-28\mu},\\
 T_\mu(c)&=c^{49}d^{147+49\mu}e^{147},\\
 T_\mu(d)&=d^{343},\qquad T_\mu(e)=e^{343}.
 \end{aligned}
 \tag{E78.6}
\]
These are the actual based operators of $\tau,\tau'$.
Put $T_-=T_{-5/2}$ and $T_+=T_{1/2}$.
Apply the source construction of E.77 with the
explicitly transported marking (E78.4). It gives
actual generators $t_s,t'_s,r_s,r'_s$ of
$D=G_k(3)$, mapping to $\tau,\tau',1,1$, with
\[
 D=
 \left\langle t_s,t'_s,r_s,r'_s\ \middle|\
 f(t_s,r_s)f(t'_s,r'_s)^{-1}=1\right\rangle_{\rm pro\text{-}3},
 \qquad f(t,r)=trt^{-1}r^{-7}.
 \tag{E78.7}
\]
All choices are made in this actual source over its
fixed quotient $D\to P$.

For $v_\mu(u,a_0)=uT_\mu(a_0)u^{-1}a_0^{-7}$,
the full local section set $H^1(k,C)$ has the
following exact compact coordinates:
\[
 \begin{gathered}
 \{(u,u',a_0,a'_0)\in C^4:
                v_{-5/2}(u,a_0)=v_{1/2}(u',a'_0)\}/C,\\
 w:(u,u',a_0,a'_0)\longmapsto
 (wuT_-(w)^{-1},\,wu'T_+(w)^{-1},\,
             wa_0w^{-1},\,wa'_0w^{-1}).
 \end{gathered}
 \tag{E78.8}
\]
The four entries are the values on
$t_s,t'_s,r_s,r'_s$, respectively.

Both lower central lifting images are evaluated as
follows. A $C_2$-valued section is represented by
\[
 \begin{gathered}
 a_0=x^A y^B c^n,\qquad u=x^U y^V c^N,\\
 a'_0=x^{A'}y^{B'}c^{n'},\qquad
 u'=x^{U'}y^{V'}c^{N'},\\
 L_\mu=3AB-AV+BU+B\mu+6n.
 \end{gathered}
 \tag{E78.9}
\]
Its exact relation is $L_{-5/2}=L'_{1/2}$, where a
prime means substitution of all six primed variables.
For the same words with zero $d,e$ coordinates put
\[
 \begin{aligned}
 R_\mu=\frac72\big(&29A^2B-7A^2V+14ABU+14AB\mu-3AB\\
 &-2AN-2AUV+AV+90An+BU^2+2BU\mu-BU\\
 &+B\mu^2-B\mu+14Un+14\mu n+42n\big),
 \end{aligned}
 \tag{E78.10}
\]
and
\[
 \begin{aligned}
 S_\mu=\frac72\big(&13AB^2-3AB-AV^2+AV+7B^2U+7B^2\mu\\
 &-2BN-BU+2BV\mu-B\mu+90Bn+4B+14Vn+42n\big).
 \end{aligned}
 \tag{E78.11}
\]
The image of $H^1(k,C)\to H^1(k,C_2)$ is exactly
\[
 \Delta R=R_{-5/2}-R'_{1/2}\in3\mathbb Z_3,\qquad
 \Delta S=S_{-5/2}-S'_{1/2}\in3\mathbb Z_3.
 \tag{E78.12}
\]
A particular lift keeps $u,u',a'_0$ unchanged and
replaces
\[
           a_0\longmapsto
              a_0d^{-\Delta R/336}e^{-\Delta S/336}.
 \tag{E78.13}
\]
Every nonempty fibre, modulo full geometric conjugacy,
has the explicit compact coordinates
\[
                  \mathbb Z_3^4\times(\mathbb Z/9)^2.
 \tag{E78.14}
\]
More precisely, after the particular lift write
$[z]=d^{z_1}e^{z_2}$. Every other lift has the
following changes, and its class has the indicated
unique coordinates:
\[
 \begin{gathered}
 a_0\mapsto a_0[z],\qquad a'_0\mapsto a'_0[z],
 \qquad u\mapsto u[b],\qquad u'\mapsto u'[b'],\\
 (z,\ b-b',\ b'\bmod342\mathbb Z_3^2).
 \end{gathered}
 \tag{E78.15}
\]
For the inverse choose the two residues of $b'$ in
$\{0,\ldots,8\}$ and add $b-b'$ to obtain $b$.
The exact remaining conjugation is
$(b,b')\mapsto(b-342w,b'-342w)$.

At the preceding stage, every abelian section has four
arbitrary vector values on the same source generators.
For their coordinates in (E78.9) define
\[
 \begin{aligned}
 D_2={}&3(AB-A'B')-\tfrac52B-\tfrac12B'\\
       &+UB-VA-U'B'+V'A'.
 \end{aligned}
 \tag{E78.16}
\]
Its class is in the image of
$H^1(k,C_2)\to H^1(k,C_1)$ exactly when
$D_2\in3\mathbb Z_3$.
A particular lift sets $n=-D_2/6$, $n'=N=N'=0$.
Every nonempty fibre has coordinates
$\mathbb Z_3^2\times\mathbb Z/3$, given by the scalar
version of (E78.15) with $c$ in place of $(d,e)$
and $48$ in place of $342$.
The abelian section group itself has the explicit
coordinates
\[
 H^1(k,C_1)\simeq\mathbb Z_3^6\times(\mathbb Z/3)^2:
 \quad(u-u',a_0,a'_0,\ u'\bmod6\mathbb Z_3^2).
 \tag{E78.17}
\]
All inverse maps retain the actual conjugating element.

For an explicit family take, in $C_2$,
\[
        a_0=y,\qquad u=x^{5/2}c^N,\qquad
        a'_0=u'=1.
 \tag{E78.18}
\]
Every member satisfies the source relation at class two.
Its exact class-three error is $e^{7(2-N)}$.
Its fibre is nonempty exactly for $N\equiv2\pmod3$;
one lift sends $r_s$ to $y e^{(N-2)/48}$.
Thus $N=0$ has empty fibre and $N=2$ already lifts
with the same displayed words. They have the same
abelian coordinates.

Every member of (E78.18) has an actual intrinsic
$N/\gamma_3N$ preimage. The proof constructs them
from a rational cusp path and two chosen geometric
lifts of $x,y$. The preimage constructed for $N=2$
comes from a full $N$-valued local section.
The preimage constructed for $N=0$ has no
$N/\gamma_4N$-valued lift. The two constructed
intrinsic sections have the same intrinsic
abelian coordinates.

The arithmetic comparison with the cubic is the
restriction to $G_k$ followed by the actual maps
\[
 \begin{gathered}
 N/\gamma_4N\twoheadrightarrow C,\qquad
 N/\gamma_3N\twoheadrightarrow C_2,\\
 \ker(N/\gamma_jN\to F/\gamma_jF)
       =h_*^{-1}(\gamma_jF)/\gamma_jN,\qquad j=3,4.
 \end{gathered}
 \tag{E78.19}
\]
It gives a commutative diagram of section sets and
the central map
$\operatorname{Lie}_3(N^{\rm ab})\to
  \operatorname{Lie}_3(F^{\rm ab})=\mathbb Z_3(3)^2$.
Consequently (E78.12) is an evaluated component of
the intrinsic wild obstruction. In particular every
constructed preimage of (E78.18) with
$N\not\equiv2\pmod3$ has empty intrinsic
class-three fibre. For an arbitrary quotient section
the kernel in (E78.19) is the retained intrinsic
lifting problem.

\emph{Manuscript: \S51.16, folios733v--736v,
scans172--178, formulas(371)--(379), Part E(c),(d)
and follow-up(4). The actual operators, source
marking and tangent transport in this quotient
are all evaluated.}
\end{statement}

\begin{proof}
\emph{Coherent roots and the literal source transport.}
The identity $\lambda^2=-3\zeta_3$ gives (E78.2):
of the two square roots of
$-3_n\zeta_{3^{n+1}}$, exactly one has $3^n$th
power $\lambda$. The same uniqueness makes its
cube equal to the preceding root. Hence these are
coherent roots, rather than separate choices at
finite levels. Since the coherent roots of $-1$
can all be taken equal to $-1$, their Kummer
character is zero. Taking Galois ratios in (E78.2)
therefore gives
\[
 2\kappa_\lambda=\kappa+\frac{\chi-1}{3},
 \qquad
 \kappa_a=\frac{\chi-1}{3}-3\kappa_\lambda,
 \tag{E78.20}
\]
which proves (E78.3).
All these identities hold on the actual $G_K$
and after restriction to $G_k$.

For the E.37 pair, the cubic residue of three
at the prime over seven with $\zeta_3\bmod7=2$
gives $t_0\equiv1\pmod3$. Positive real roots
give $\kappa(c_{\rm ar})=0$, so
$\kappa(\tau'_0)=-t_0$.
The cocycle law now gives the exact identities
\[
 \begin{gathered}
 \chi(\rho_0)=1,\quad\kappa(\rho_0)=-2t_0,\quad
 \kappa(h_0)=(t_0-1)/6,\\
 \kappa(\tau)=1,\quad\kappa(\tau')=-1,\quad
 \kappa_a(h_0)=(1-t_0)/4.
 \end{gathered}
 \tag{E78.21}
\]
The exponent in (E78.4) is integral because
$t_0-1\in3\mathbb Z_3$. The endomorphism of the
free group $P$ sending its original basis to
$(\tau,\tau')$ induces the identity on its
Frattini quotient. It is consequently a continuous
automorphism. It commutes with $c_{\rm ar}$ and
preserves $\chi$. Formula (E78.4), together with
the inverse of this actual free-group
automorphism, is the source marking map.

This transport also has an exact geometric
description. Before it, the two operators of
E.54 are
$\alpha_{7,-1-3t_0/2}$ and
$\alpha_{7,-1+3t_0/2}$.
The operator of $h_0$ is
$\alpha_{1,(1-t_0)/4}
 =\operatorname{Ad}(x^{(1-t_0)/4})$.
The composition law in E.54 conjugates the first
operator to $\alpha_{7,-5/2}$.
For the second one, $c_{\rm ar}h_0c_{\rm ar}$
has the opposite Kummer coordinate and conjugates
it to $\alpha_{7,1/2}$. These are literal
equalities of based automorphisms of $C$.
Thus the change of Frobenius prolongations and
the induced inner path operators have both been
specified.

For completeness, implement the new marking in
the source construction as follows. In E.75's
initial splitting replace the horizontal lifts
by $j_0(h_0)t_{s,0}j_0(h_0)^{-1}$ and its
$c_{\rm ar}$ conjugate, leaving the initial
vertical pair unchanged. In the free relative
presentation these substitutions form a
Frattini-invertible automorphism and induce
exactly (E78.4) on $P$.
Run the oriented Fox change and the convergent
relative corrections of E.75--E.77 with these
horizontal generators. Their proofs use the
actual free basis, its swapped involution and
its two cyclotomic values seven, all retained
by (E78.4). Every later correction is in the
kernel of the same map to $P$.
This gives (E78.7) in the actual local group.
The change from the old source splitting to
the new one is the actual cocycle
$j_{\rm new}(g)j_{\rm old}(g)^{-1}$ in that
kernel, with the change-of-section formula
of (E77.2) and (E77.19).
In particular no equality of the original
and transported Frobenius lifts is used.

\emph{The genuine operators and the source equation.}
E.54 gives
\[
\alpha_{\chi,k}(y)=
 x^k c^{-(\chi^2-1)/24}y^\chi
 c^{(\chi^2-1)/24}x^{-k}.
\]
Its degree-two coefficient follows by integral
collection in the genuine three-cycle relation,
with the positive commutator convention checked
there; the original relation and coefficient are
also in \cite[Theorem2, equation(5), \S3 and
Corollary12]{A_Furusho2010Hexagon}.
The higher part of the conjugator is central
in $C$ and acts trivially by conjugation there.
Substitution of (E78.5) in E.54's full Hall
images gives (E78.6), including its
$c,d,e$ terms.

The action on $C$ factors through $P$:
the two characters in (E78.3) are unramified
outside three, and their action is E.54's
explicit pro-three image. For a continuous
$C$-valued cocycle, its graph in the
semidirect product with that pro-three image
is a homomorphism into a pro-three group.
It therefore factors through $D=G_k(3)$.
Assigning arbitrary four values to the free
generators of the presentation (E78.7)
defines a homomorphism into $C\rtimes P$.
The relator maps precisely to
$v_{-5/2}(u,a_0)v_{1/2}(u',a'_0)^{-1}$.
It vanishes exactly under (E78.8), proving
both existence and exhaustiveness of the
section coordinates. Geometric conjugation
gives the four formulas displayed there.

\emph{Integral collection and both images.}
The Hall law from E.54 is
\[
 \begin{aligned}
 &(a,b,n,r,s)(A,B,N,R,S)\\
 &\quad=\bigl(a+A,\ b+B,\ n+N-Ab,\\
 &\qquad r+R+bA(A+1)/2-An,\\
 &\qquad s+S+Ab(b+1)/2-Bn+AbB\bigr).
 \end{aligned}
 \tag{E78.22}
\]
Every denominator is a unit at three. It is an
identity for integral powers by collection,
and hence for all five three-adic coordinates
by density and continuity. In these coordinates
the actual operator has the complete expression
\[
 \begin{aligned}
 T_\mu(A,B,n,r,s)=\bigl(&7A,\ 7B,\ 7\mu B+49n,\\
 &7\mu(\mu-1)B/2+(147+49\mu)n+343r,\\
 &(14-28\mu)B-49\mu B(B-1)/2+147n+343s\bigr).
 \end{aligned}
 \tag{E78.23}
\]
Substitute (E78.22)--(E78.23) into the four
factors defining $v_\mu$. Collection yields
the exact identity
\[
        v_\mu(u,a_0)=c^{7L_\mu}d^{R_\mu}e^{S_\mu}
 \tag{E78.24}
\]
for the zero-central Hall lifts in (E78.9).
The explicit polynomials are (E78.10)--(E78.11).
They may alternatively be verified coefficient
by coefficient in the truncated integral
Magnus algebra; that independent computation
is part of the certificate below.

In $C_2$ the relator is therefore
$7(L_{-5/2}-L'_{1/2})$.
Its part independent of $n,n'$ is $7D_2$,
and its remaining part is $42(n-n')$.
This proves (E78.16) and its actual particular
lift. At class three the common $c$ coordinate
cancels: the subgroup generated by $c,d,e$
is abelian, so the exact remaining error is
$d^{\Delta R}e^{\Delta S}$.
Changing the central coordinates of $a_0,a'_0$
by $p,p'\in\mathbb Z_3^2$ changes this error by
\[
                   (343-7)(p-p')=336(p-p').
 \tag{E78.25}
\]
Changing the central coordinates of $u,u'$
does not change it. Since
$336\mathbb Z_3=3\mathbb Z_3$, this proves
the necessity and sufficiency in (E78.12)
and the particular lift (E78.13).

\emph{The full conjugacy fibres and their topology.}
At either step let the central kernel be $A_j$,
where $A_2=\mathbb Z_3(2)c$ and
$A_3=\mathbb Z_3(3)d\oplus\mathbb Z_3(3)e$.
After one solution, the differences on
$r_s,r'_s$ satisfy
$(7^j-7)(p-p')=0$. These lattices are
torsion-free and $j=2,3$, so $p=p'$.
The differences on $t_s,t'_s$ are arbitrary
$b,b'$, modulo exactly
\[
       (b,b')\longmapsto
             (b-(7^j-1)w,\ b'-(7^j-1)w).
 \tag{E78.26}
\]
This proves the stated coordinates and their
inverse with $7^2-1=48$ and $7^3-1=342$.
More explicitly the normalizing conjugator is
$c^{(b'-\epsilon)/48}$ at class two and
$[(b'-\epsilon)/342]$ at class three, where
$\epsilon$ is the chosen finite residue.

There is no additional quotient from a lower
stabilizer. Indeed, a nonidentity element fixed
by a lower twisted section has a first
nonzero Hall graded coordinate. The operator
$\operatorname{Ad}(u)T_-$ on degree $j$
is multiplication by $7^j$. The equation
$(7^j-1)v=0$ has no nonzero solution in the
torsion-free graded lattice. Thus the lower
fixed subgroup is trivial.
Every geometric conjugacy of lower sections
consequently has a unique lower conjugator.
Lifting it gives precisely the central
freedom already divided out in (E78.26).
At the abelian stage the same calculation
uses $7-1=6$ and gives (E78.17) directly.

All coordinate spaces and the relator-zero
sets are compact. Their group actions are
continuous and their orbit spaces are
Hausdorff. The orbit map for a fixed lower
section has trivial stabilizer, hence is a
homeomorphism from the compact lower group
to that orbit. Division by $48$ or $342$
on its image is continuous, and selection
of the indicated finite residues is
continuous. These facts prove the asserted
homeomorphisms, including the inverses
with their actual normalizing conjugators.

Substituting (E78.18) in (E78.24)
gives $L=R=0$ and $S=7(2-N)$.
Equation (E78.25) proves every assertion
about its quotient fibres. The maps in (E78.19)
and their kernels are induced by the
actual surjection $h_*$ of E.54.
Its path comparison is the group
isomorphism (E54.17), with the specific
roots (E78.2), so it also intertwines the
arithmetic semidirect products and all
section and central obstruction maps.

\emph{Actual intrinsic preimages and a wild nonlifting
section.}
For the full based action used here, the rational coordinate
isomorphism $[x:y:z]\mapsto[z:y:x]$ sends E.34's cusp
$[0:-1:1]$ to $[1:-1:0]$. It sends the source parameter
$s_0=-x/z$ to the parameter $s=z/x$ with $s=-s_0$.
Choose the compatible $3^n$-th roots of $-1$ all equal
to $-1$. These roots are rational and compatible, so
the resulting punctured-disc comparison is Galois invariant.
The coordinate isomorphism followed by this tangential
comparison therefore transports E.34's genuine based
$N$-action to the present one. In particular it still
factors through the same $P$ over $K$, as required when
using the actual source generators in (E78.7).

Choose any of the three $K$-rational cusps $B_i$
where $1-h$ has order one, as in (E54.16).
An actual geometric path from the original tangent
to a positive tangent at this cusp gives an inertia
generator $d_*\in N$ and an element $u_*\in N$
with
\[
                 u_*T(d_*)u_*^{-1}=d_*^7 .
 \tag{E78.27}
\]
Here $T$ is the genuine operator of $\tau$ on $N$.
Explicitly, for the path $p$ the element
$p\,\tau(p)^{-1}$ is such a conjugator, because
the local cusp action raises its positive
inertia generator to the seventh power.
The image of $d_*$ under $h_*$ is conjugate to
$y$, with exponent one. If it is $g y g^{-1}$,
choose $v\in N$ with $h_*(v)=g^{-1}$ and replace
\[
 d_*\longmapsto vd_*v^{-1},\qquad
 u_*\longmapsto vu_*T(v)^{-1}.
 \tag{E78.28}
\]
Surjectivity of the actual $h_*$ supplies $v$,
and (E78.28) is the exact morphism preserving
(E78.27). We now have $h_*(d_*)=y$.

Write the image of $u_*$ in $C$ in Hall coordinates.
The identity
$v_{-5/2}(h_*(u_*),y)=1$ forces its
$x$ coordinate to be $5/2$. Equations
(E78.10)--(E78.11) then force its $c$
coordinate to be $2-5V/2$, where $V$ is
its $y$ coordinate; the $d,e$ coordinates
are retained without evaluation.
Replace the actual generator value $u_*$
by $d_*^{-V}u_*$. Left multiplication by
this cusp power preserves (E78.27).
The exact Hall law gives
\[
 \begin{gathered}
 h_*(d_*^{-V}u_*)\bmod\gamma_3F=x^{5/2}c^2,\\
 (u,u',a_0,a'_0)=(d_*^{-V}u_*,\,1,\,d_*,\,1).
 \end{gathered}
 \tag{E78.29}
\]
These four actual $N$-values satisfy both sides
of the relator in (E78.7) separately with
value one. They therefore define a full
$N$-valued local section. Formula (E78.29)
specifies a choice of section values, rather
than identifying distinct original sections
without a map.

Choose actual $x_*,y_*\in N$ with images $x,y$
and put $w_*=[x_*,y_*]$. The class of $w_*$
is central in $N/\gamma_3N$ and its image
is $c$. In that quotient replace only the
$t_s$ value in (E78.29) by
\[
                 d_*^{-V}u_*\,w_*^{\,N-2}.
 \tag{E78.30}
\]
The relator remains exactly one: its changed
factor conjugates $T(d_*)$ by a central
element. Thus (E78.30) constructs an actual
intrinsic class-two section for every
$N\in\mathbb Z_3$, with image (E78.18).
All their intrinsic abelian values agree.
At $N=2$ the original full section is an
intrinsic lift. At $N=0$, the nonzero
class-three obstruction already computed
in its actual quotient excludes every
intrinsic class-three lift.
This proves both intrinsic assertions by
explicit source and geometric maps.

The self-contained program
\texttt{a\_wild\_class3.py} proves the full
symbolic identities for $L_\mu,R_\mu,S_\mu$
and for the central change $336$. It checks
640 identities in an independent associative
Magnus model, all five Hall images for both
actual operators, eight exact source/root
transports, 67 class-two lifting cases,
seven class-three lifting cases, 180 abelian
and 201 class-two fibre inverses, and 35
complete central-fibre coordinates and
inverses. Its symbolic family calculation
retains the variable $N$, and its symbolic
cusp-path normalization retains $V$ and
both unevaluated central path coordinates.
All final checks
pass; the all-input conclusions follow
from the polynomial identities and the
proved actual source presentation.
\end{proof}
